% =====================================================================
%  Дополнения preamble.tex для отчётов по лабораторным работам.
%  Подключается ПОСЛЕ \input{preamble} в lab/lab_report.tex.
%
%  Компиляция требует ключ -shell-escape (нужен пакету minted для
%  подсветки синтаксиса кода через Pygments). См. lab/README.md.
% =====================================================================

\usepackage{float} % на случай отдельной компиляции без полной цепочки minted

% ---------------------------------------------------------------------
% Сквозная нумерация рисунков/таблиц/листингов/формул (1, 2, 3, ...)
% вместо "привязанной к разделу" (preamble.tex рассчитан на курсовые/
% дипломные проекты с нумерованными \section; в лабораторных работах
% используются ненумерованные заголовки \labpart, поэтому нумерацию
% по разделам здесь отключаем).
% ---------------------------------------------------------------------
\AtBeginDocument{%
  \renewcommand{\thefigure}{\arabic{figure}}
  \renewcommand{\thetable}{\arabic{table}}
  \renewcommand{\thelisting}{\arabic{listing}}
  \renewcommand{\theequation}{\arabic{equation}}
}

% ---------------------------------------------------------------------
% Код-листинги (пакет minted)
% ---------------------------------------------------------------------

\usemintedstyle{friendly}
\setminted{
  fontsize=\small,
  breaklines=true,
  breakanywhere=true,
  linenos=true,
  framesep=2mm,
  tabsize=4,
}

% Листинг из отдельного файла с исходным кодом.
% \codelistingfile[опции minted]{язык}{путь/к/файлу}{Подпись листинга}{метка-без-префикса}
% Пример: \codelistingfile{python}{lab/code/example.py}{Пример скрипта}{example-script}
%         -> ссылка в тексте: \autoref{lst:example-script}
\newcommand{\codelistingfile}[5][]{%
  \begin{listing}[H]
    \inputminted[#1]{#2}{#3}
    \caption{#4}
    \label{lst:#5}
  \end{listing}%
}

% Листинг, код которого пишется прямо в тексте отчёта.
% \begin{codelisting}[опции minted]{язык}{Подпись листинга}{метка-без-префикса}
%   ...код...
% \end{codelisting}
\newenvironment{codelisting}[4][]{%
  \def\codelistingCurrentCaption{#3}%
  \def\codelistingCurrentLabel{#4}%
  \VerbatimEnvironment
  \begin{listing}[H]
  \begin{minted}[#1]{#2}%
}{%
  \end{minted}
  \caption{\codelistingCurrentCaption}
  \label{lst:\codelistingCurrentLabel}
  \end{listing}%
}

% ---------------------------------------------------------------------
% Диаграммы PlantUML (см. lab/diagrams/README.md)
% ---------------------------------------------------------------------

% \pumldiagram{имя-файла-без-расширения}{Подпись}{метка-без-префикса}{доля \textwidth}
% Пример: \pumldiagram{sequence-query}{Диаграмма последовательности обработки запроса}{sequence-query}{0.9}
%         -> ожидает файл lab/diagrams/generated/sequence-query.png
%         -> ссылка в тексте: \autoref{fig:sequence-query}
\newcommand{\pumldiagram}[4]{%
  \begin{figure}[H]
    \centering
    \includegraphics[width=#4\textwidth]{lab/diagrams/generated/#1}
    \caption{#2}
    \label{fig:#3}
  \end{figure}%
}

% ---------------------------------------------------------------------
% Заголовки тематических блоков лабораторной работы
% ---------------------------------------------------------------------
% Ненумерованный, жирный, отцентрированный заголовок без изменения
% регистра (как в примерах отчётов по лабораторным). По умолчанию НЕ
% начинает новую страницу — если нужен разрыв, поставьте \clearpage
% (или \newpage) перед вызовом сами.
%
% \labpart{Название блока}
\newcommand{\labpart}[1]{%
  \par\vspace{1.5em}%
  \phantomsection
  \begin{center}\bfseries #1\end{center}
  \addcontentsline{toc}{section}{#1}
  \markboth{#1}{}%
}

% Та же команда, но принудительно с новой страницы.
% \labpartbreak{Название блока}
\newcommand{\labpartbreak}[1]{\clearpage\labpart{#1}}
